\section{The deep large sieve for Fermat-quotient characters}\label{sec:dls}

This section is self-contained and fully proved; all identities and bounds
have additionally been verified numerically (Appendix~\ref{app:numerics}).

\begin{definition}[Fermat quotient; depth-one characters]\label{def:fq}
Let $q$ be an odd prime. For $(n,q)=1$ the \emph{Fermat quotient} is
\[
\fq(n)\;\equiv\;\frac{n^{q-1}-1}{q}\pmod q .
\]
It depends only on $n\bmod q^2$ and defines a surjective homomorphism
$(\Z/q^2)^*\to(\Z/q,+)$:
\begin{equation}\label{eq:fq-hom}
\fq(mn)=\fq(m)+\fq(n),
\qquad
\fq(1+qu)\equiv-u\pmod q .
\end{equation}
Its kernel is the \emph{Teichm\"uller subgroup} $T_q$ (the $(q-1)$-th
roots of unity mod $q^2$; $|T_q|=q-1$). The \emph{depth-one characters}
are
\[
\chi_\lambda(n)\;:=\;\e{\frac{\lambda\,\fq(n)}{q}},\qquad
\lambda\in\Z/q;
\]
they are exactly the characters of $(\Z/q^2)^*$ trivial on $T_q$, and for
$\lambda\ne0$ the conductor is $q^2$.
\end{definition}

\begin{proof}[Proof of \eqref{eq:fq-hom}]
Multiplicativity: $(mn)^{q-1}-1=(m^{q-1}-1)n^{q-1}+(n^{q-1}-1)$, and
$n^{q-1}\equiv1\pmod q$, so dividing by $q$ and reducing mod $q$ gives the
claim. Second identity:
$(1+qu)^{q-1}=1+(q-1)qu+\binom{q-1}{2}q^2u^2+\cdots\equiv1-qu\pmod{q^2}$,
so $\fq(1+qu)\equiv-u$. Surjectivity follows from the second identity
applied to $u=0,\ldots,q-1$; the kernel has index $q$, order $q-1$, and
consists of the solutions of $y^{q-1}\equiv1\pmod{q^2}$, which form the
Teichm\"uller subgroup. For $\lambda\ne0$, $\chi_\lambda$ is nontrivial on
$1+q\Z/q^2$ by the second identity, hence has conductor $q^2$.
\end{proof}

\begin{remark}[Wieferich orthogonality]
Summing the characters,
\[
\sum_{\lambda\bmod q}\chi_\lambda(y)
=q\cdot\mathbf 1\big[y^{\,q-1}\equiv1\ (\mathrm{mod}\ q^2)\big]:
\]
orthogonality over the deep family is a Wieferich-type condition ---
pointwise intractable, but average-friendly. This is the algebraic engine
of everything below.
\end{remark}

\begin{proposition}[fiber bound]\label{prop:fiber}
\statusPM\quad
For any $N\ge1$ and $v\in\Z/q$,
\[
\#\{n\le N:\ (n,q)=1,\ \fq(n)=v\}
\;\le\;(q-1)\Big(\fl{\frac{N}{q^2}}+1\Big)
\;\le\;\frac Nq+q .
\]
\end{proposition}

\begin{proof}
$\fq$ is constant on residue classes mod $q^2$; exactly $q-1$ such classes
(one coset of $T_q$) take the value $v$; each class meets $[1,N]$ in at
most $\fl{N/q^2}+1$ points.
\end{proof}

\begin{proposition}[Teichm\"uller energy / Wieferich-pair count]\label{prop:energy}
\statusPM\quad
\[
E(N;q):=\#\Big\{(n_1,n_2)\in[1,N]^2:\ (n_1n_2,q)=1,\
n_1^{q-1}\equiv n_2^{q-1}\ (\mathrm{mod}\ q^2)\Big\}
\;\le\;\frac{N^2}{q}+qN .
\]
\end{proposition}

\begin{proof}
The condition is $\fq(n_1)=\fq(n_2)$, so $E=\sum_vf_v^2$ with $f_v$ the
fiber counts; then
$E\le(\max_vf_v)\sum_vf_v\le(N/q+q)\,N$ by
Proposition~\ref{prop:fiber}.
\end{proof}

\begin{theorem}[deep large sieve, single modulus]\label{thm:dls}
\statusPM\quad
For any complex weights $(a_n)_{n\le N}$ supported on $(n,q)=1$,
\[
\sum_{\lambda\bmod q}\Big|\sum_{n\le N}a_n\chi_\lambda(n)\Big|^2
\;\le\;\big(N+q^2\big)\,\|a\|_2^2 .
\]
\end{theorem}

\begin{proof}
By orthogonality of the additive characters of $\Z/q$, the left side
equals $q\sum_{v}|A_v|^2$ with $A_v=\sum_{\fq(n)=v}a_n$. Cauchy--Schwarz
within each fiber gives $|A_v|^2\le f_v\sum_{\fq(n)=v}|a_n|^2$, so by
Proposition~\ref{prop:fiber}
$\sum_v|A_v|^2\le(N/q+q)\|a\|_2^2$.
\end{proof}

\begin{remark}
The family has $q$ characters of conductor $q^2$; $(N+q^2)$ is the natural
conductor-limited analogue of the classical $(N+Q^2)$ and is sharp for
$N\ge q^2$.
\end{remark}

\begin{corollary}[per-character RMS; nontriviality range]\label{cor:rms}
\statusPM\quad
For $a$ the indicator of an interval of length $N\le q^2$,
\[
\Big(\frac{1}{q-1}\sum_{\lambda\ne0}\Big|\sum_{n}a_n\chi_\lambda(n)\Big|^2\Big)^{1/2}
\;\ll\;\sqrt{qN}\,,
\]
a saving of $\sqrt{N/q}$ over the trivial bound $N$: nontrivial whenever
$N>q^{1+\varepsilon}$. (For $N>q^2$ the RMS bound is $N/\sqrt q$.)
\end{corollary}

\begin{corollary}[crude average over moduli]\label{cor:crude}
\statusP\quad
Summing Theorem~\ref{thm:dls} over primes $q\sim Q$: the total mean square
over the family $\{(q,\lambda):\lambda\ne0\}$ (of size $\asymp Q^2/\log Q$)
is $\ll(N^2+Q^3N)/\log Q$, i.e.\ per-character RMS $\ll N/Q+\sqrt{QN}$.
This is no cross-moduli gain --- it merely sums the single-$q$ bounds. The
genuine strengthening (the program's flagged main theoretical question) is
a tailored $(q,\lambda)$-averaged sieve exploiting cross-moduli
Wieferich-pair counting: for fixed $n_1\ne n_2$ the primes $q$ with
$q^2\mid n_1^{q-1}-n_2^{q-1}$ should be far sparser than fixed-$q$ counting
sees.
\end{corollary}

\begin{corollary}[progression-restricted version]\label{cor:ap}
\statusPM\quad
Fix $c$ with $(c,q)=1$ and let $a$ be supported on
$n\equiv c\ (\mathrm{mod}\ q)$, $n\le N$. Then
\[
\sum_{\lambda\bmod q}\Big|\sum_na_n\chi_\lambda(n)\Big|^2
\;\le\;\Big(\frac Nq+q\Big)\|a\|_2^2 .
\]
\end{corollary}

\begin{proof}
On the progression $n=n_0+qs$ (any fixed $n_0\equiv c$ with
$(n_0,q)=1$) one has $n\equiv n_0(1+qs\bar n_0)\pmod{q^2}$ with
$\bar n_0n_0\equiv1\ (q)$, so by \eqref{eq:fq-hom}
$\fq(n)=\fq(n_0)-s\bar n_0\bmod q$: as $s$ runs over a complete residue
system mod $q$, $\fq(n)$ visits each value exactly once. Hence each
$\fq$-fiber meets the progression in a \emph{single} residue class mod
$q^2$, the fiber count improves to $N/q^2+1$, and the proof of
Theorem~\ref{thm:dls} gives the stated bound. Equivalently, this is the
classical additive large sieve in the coordinate $s=(n-n_0)/q$.
\end{proof}

\begin{remark}[parameter conclusions for Lemma D]
With band parameters $p=x^u$, $q=x^{u'}$, $u,u'\in(1/3,1/2)$: the
$a$-side (length $N_a=x^{1-u}$) always lies in the nontrivial range
($u'<\tfrac12<1-u$), with unconditional saving $x^{(1-u-u')/2}$ after
Cauchy--Schwarz in $p$; the $p$-side (length $x^u$) is nontrivial only
for $u>u'$, half the band. These are mean bounds over $\lambda$; the
$\lambda$-aggregation is the subject of
Proposition~\ref{prop:largelambda}.
\end{remark}
